\section{Convergence and Complexity Analysis}
\label{sec:complexity}

In this section, we establish the convergence properties and oracle complexity of TOBYQA under fully-linear models.

\subsection{Standing Assumptions and Stochastic Framework}
\label{sec:assumptions}

We analyze TOBYQA under the classical framework of fully-linear models~\cite{ConnScheinbergVicente2009,ChenMenickellyScheinberg2018,CartisRoberts2019}. Let $\mathcal{F}_{k-1} := \sigma(\vct{x}_0, \dots, \vct{x}_k, \mathcal{Y}_0, \dots, \mathcal{Y}_k)$ denote the filtration representing the history of the algorithm up to the beginning of iteration $k$. We impose the following standard assumptions on the latent objective function and the observation channel.

\begin{assumption}
\label{ass:smoothness}
The objective function $f: \mathbb{R}^n \to \mathbb{R}$ is continuously differentiable with $L_g$-Lipschitz continuous gradient:
\begin{equation}\label{eq:grad-lipschitz}
  \|\nabla f(\vct{x}) - \nabla f(\vct{y})\|_2 \;\le\; L_g\,\|\vct{x} - \vct{y}\|_2, \qquad \forall\, \vct{x}, \vct{y} \in \mathbb{R}^n,
\end{equation}
and is bounded from below by $f_{\mathrm{low}} > -\infty$ on $\mathbb{R}^n$.
\end{assumption}

\begin{assumption}
\label{ass:radius-bounds}
\textup{\textbf{(Active-set radius bounds.)}}
There exist constants $0<\delta_{\mathrm{lo}}\le\delta_{\mathrm{hi}}<\infty$ such that, at every iteration prior to termination,
\begin{equation}\label{eq:radius-bounds}
  \delta_{\mathrm{lo}}\;\le\;\delta_k\;\le\;\delta_{\mathrm{hi}},
  \qquad
  \delta_k:=\max_{i=1,\dots,m}\|\vct{x}_i-\vct{x}_k\|_2.
\end{equation}
Such active-set radius bounds are standard in derivative-free complexity analysis~\cite{ConnScheinbergVicente2009,CartisRoberts2019}.
\end{assumption}

\begin{assumption}
\label{ass:fully-linear}
\textup{\textbf{(Fully-linear gradient accuracy.)}}
There exists a constant $\kappa_{eg} > 0$ such that, at every iteration $k \ge 0$, the recovered model gradient $\vct{g}_k$ satisfies
\begin{equation}\label{eq:grad-error-bound}
  \|\vct{g}_k - \nabla f(\vct{x}_k)\|_2 \;\le\; \kappa_{eg}\,\zeta_k,
\end{equation}
where $\zeta_k := \max\{\delta_k, \sigma_\varepsilon / \delta_k\}$ and $\delta_k$ is the active sample radius in~\eqref{eq:radius-bounds}.
\end{assumption}

Assumption~\ref{ass:fully-linear} is the standard accuracy condition for fully-linear models in derivative-free optimization~\cite{ConnScheinbergVicente2009}. Its error scale combines the sample radius $\delta_k$, which controls spatial truncation, and the ratio $\sigma_\varepsilon / \delta_k$, which accounts for noise amplification.

For the observation channel, Proposition~\ref{prop:ared-decomp} and Theorem~\ref{thm:exact-decoupling} show that the affine drift rate is represented by the temporal coefficient and does not enter the recovered spatial gradient. They do not imply that the drift-free time coefficient $\gamma_k^{\mathrm{stat}}$ vanishes. Sequential space--time correlation can produce a nonzero $\gamma_k^{\mathrm{stat}}$ even without observation noise. We therefore define the conditional acceptance bias
\begin{equation}\label{eq:ared-expectation}
  b_k \;:=\; \mathbb{E}[\mathrm{Ared}_k \mid \mathcal{F}_{k-1}]
  - \bigl(f(\vct{x}_k)-f(\vct{x}_k+\vct{s}_k)\bigr).
\end{equation}
Under affine drift, relation~\eqref{eq:ared-affine-residual} identifies the deterministic space--time aliasing contribution as $\gamma_k^{\mathrm{stat}}\Delta t_k$; under $C^2$ drift, $b_k$ also contains the nonlinear temporal truncation error. The following assumption controls this conditional bias and the remaining stochastic innovation.

\begin{assumption}
\label{ass:drift-regime}
\textup{\textbf{(Acceptance-channel residual.)}}
There exists $\varepsilon > 0$ such that, whenever the incumbent satisfies $\|\nabla f(\vct{x}_k)\|_2 \ge \varepsilon$, the conditional bias in~\eqref{eq:ared-expectation} obeys
\begin{equation}\label{eq:residual-drift-bound}
  |b_k| \;\le\; \tfrac{1}{4}\,\mathrm{Pred}_k,
\end{equation}
and the centered innovation $\mathrm{Ared}_k-\mathbb{E}[\mathrm{Ared}_k\mid\mathcal F_{k-1}]$ is conditionally sub-Gaussian with variance proxy $2\sigma_\varepsilon^2$.
\end{assumption}

\subsection{Reliable Regularization Regime and Step Decrease}
\label{sec:reliable-regime}

Recall from Section~\ref{sec:arc-decision} that the Zero-Hessian closed-form cubic trial step $\vct{s}_k$ and the corresponding predicted reduction $\mathrm{Pred}_k$ satisfy:
\begin{equation}\label{eq:step-pred-relations}
  \vct{s}_k \;=\; -\frac{\vct{g}_k}{\sqrt{\sigma_k \|\vct{g}_k\|_2}},
  \qquad
  \mathrm{Pred}_k \;=\; \|\vct{g}_k\|_2^{3/2}\,\sigma_k^{-1/2} \;=\; \|\vct{g}_k\|_2\,\|\vct{s}_k\|_2,
\end{equation}
with $\|\vct{s}_k\|_2 = \sqrt{\|\vct{g}_k\|_2 / \sigma_k}$.
We first give a sufficient lower bound on the cubic regularization weight $\sigma_k$ for decrease in the latent objective $f$.

\begin{lemma}
\label{lem:reliable-sigma}
\textup{\textbf{(Descent under sufficient regularization.)}}
Let Assumption~\ref{ass:smoothness} hold. If the cubic regularization weight satisfies
\begin{equation}\label{eq:sigma-threshold}
  \sigma_k \;\ge\; \sigma^*(\vct{g}_k) \;:=\; \frac{4L_g^2}{\|\vct{g}_k\|_2},
\end{equation}
then the objective decrease along the candidate trial step $\vct{s}_k$ satisfies
\begin{equation}\label{eq:true-decrease-bound}
  f(\vct{x}_k) - f(\vct{x}_k + \vct{s}_k) \;\ge\; \tfrac{3}{4}\,\mathrm{Pred}_k \;-\; \|\vct{g}_k - \nabla f(\vct{x}_k)\|_2\,\|\vct{s}_k\|_2.
\end{equation}
\end{lemma}

\begin{proof}
By the classical descent lemma for $L_g$-smooth functions (Assumption~\ref{ass:smoothness}), we have
\begin{align*}
  f(\vct{x}_k + \vct{s}_k) 
  &\;\le\; f(\vct{x}_k) + \nabla f(\vct{x}_k)^\top \vct{s}_k + \tfrac{L_g}{2}\|\vct{s}_k\|_2^2 \\
  &\;=\; f(\vct{x}_k) + \vct{g}_k^\top \vct{s}_k + \tfrac{L_g}{2}\|\vct{s}_k\|_2^2 + \bigl(\nabla f(\vct{x}_k) - \vct{g}_k\bigr)^\top \vct{s}_k.
\end{align*}
Using $\vct{g}_k^\top \vct{s}_k = -\|\vct{g}_k\|_2 \|\vct{s}_k\|_2 = -\mathrm{Pred}_k$ and Cauchy--Schwarz, rearranging terms yields
\begin{equation}\label{eq:descent-intermediate}
  f(\vct{x}_k) - f(\vct{x}_k + \vct{s}_k) \;\ge\; \mathrm{Pred}_k \;-\; \tfrac{L_g}{2}\|\vct{s}_k\|_2^2 \;-\; \|\vct{g}_k - \nabla f(\vct{x}_k)\|_2\,\|\vct{s}_k\|_2.
\end{equation}
We now evaluate the ratio of the quadratic curvature penalty $\tfrac{L_g}{2}\|\vct{s}_k\|_2^2$ to the predicted decrease $\mathrm{Pred}_k$. Substituting the closed-form identities~\eqref{eq:step-pred-relations} gives:
\begin{equation}\label{eq:curvature-ratio}
  \frac{\tfrac{L_g}{2}\|\vct{s}_k\|_2^2}{\mathrm{Pred}_k}
  \;=\;
  \frac{\tfrac{L_g}{2}\bigl(\|\vct{g}_k\|_2 / \sigma_k\bigr)}{\|\vct{g}_k\|_2^{3/2} \sigma_k^{-1/2}}
  \;=\;
  \frac{L_g}{2\sqrt{\sigma_k \|\vct{g}_k\|_2}}.
\end{equation}
Under condition~\eqref{eq:sigma-threshold}, we have $\sigma_k \|\vct{g}_k\|_2 \ge 4L_g^2$, which implies $\sqrt{\sigma_k \|\vct{g}_k\|_2} \ge 2L_g$. Substituting this inequality into~\eqref{eq:curvature-ratio} yields
\[
  \frac{\tfrac{L_g}{2}\|\vct{s}_k\|_2^2}{\mathrm{Pred}_k} \;\le\; \frac{L_g}{2(2L_g)} \;=\; \tfrac{1}{4},
  \qquad\text{whence}\qquad
  \tfrac{L_g}{2}\|\vct{s}_k\|_2^2 \;\le\; \tfrac{1}{4}\,\mathrm{Pred}_k.
\]
Substituting $\tfrac{L_g}{2}\|\vct{s}_k\|_2^2 \le \tfrac{1}{4}\mathrm{Pred}_k$ into~\eqref{eq:descent-intermediate} yields~\eqref{eq:true-decrease-bound}, completing the proof.
\end{proof}

\begin{remark}
\label{rem:feedback-control}
\textup{\textbf{(Interpretation of the regularization scale.)}}
Condition~\eqref{eq:sigma-threshold} identifies the scale at which the curvature term in the descent bound is at most one quarter of the predicted reduction. In Algorithm~\ref{alg:tobyqa}, changing $\sigma_k$ changes the step length through $\|\vct{s}_k\|_2=\sqrt{\|\vct{g}_k\|_2/\sigma_k}$: rejection increases $\sigma_k$, whereas an accepted step with $r_k>\eta$ decreases it. Thus $\sigma_k$ controls the step length without a separately maintained trust-region radius.
\end{remark}

\subsection{Rejection Probability and Regularization Confinement}
\label{sec:supermartingale}

\begin{lemma}
\label{lem:reject-prob}
\textup{\textbf{(Rejection probability in the reliable regime.)}}
Let Assumption~\ref{ass:smoothness} hold, let the residual drift bound~\eqref{eq:residual-drift-bound} hold at iteration $k$, and suppose that
\begin{equation}\label{eq:reliable-conditions}
  \sigma_k \;\ge\; \gamma_2\,\sigma^*(\vct{g}_k)
  \qquad\text{and}\qquad
  \|\vct{g}_k - \nabla f(\vct{x}_k)\|_2\,\|\vct{s}_k\|_2 \;\le\; \tfrac{1}{4}\,\mathrm{Pred}_k.
\end{equation}
Then the conditional rejection probability obeys
\begin{equation}\label{eq:reject-prob}
  \mathbb{P}\bigl(\text{iteration } k \text{ is rejected} \,\big|\, \mathcal{F}_{k-1}\bigr)
  \;\le\;
  e^{-\kappa_{\mathrm{tol}}^2/8}
  \;=:\; p',
\end{equation}
and $p' \le \tfrac{1}{2}$ whenever $\kappa_{\mathrm{tol}} \ge \sqrt{8\ln 2}$.
\end{lemma}

\begin{proof}
Under $\sigma_k \ge \gamma_2\sigma^*(\vct{g}_k)$ we have $\sigma_k\|\vct{g}_k\|_2 \ge 4L_g^2$, so Lemma~\ref{lem:reliable-sigma} gives $\Delta f_k \ge \tfrac{3}{4}\mathrm{Pred}_k - \|\vct{g}_k - \nabla f(\vct{x}_k)\|_2\|\vct{s}_k\|_2 \ge \tfrac{1}{2}\mathrm{Pred}_k$. Together with~\eqref{eq:residual-drift-bound} and the expectation~\eqref{eq:ared-expectation}, this yields
\[
  \mathbb{E}[\mathrm{Ared}_k \mid \mathcal{F}_{k-1}]
  \;=\; \Delta f_k + b_k
  \;\ge\; \tfrac{1}{2}\mathrm{Pred}_k - \tfrac{1}{4}\mathrm{Pred}_k
  \;=\; \tfrac{1}{4}\mathrm{Pred}_k \;>\; 0.
\]
The deviation $\mathrm{Ared}_k - \mathbb{E}[\mathrm{Ared}_k \mid \mathcal{F}_{k-1}]$ is sub-Gaussian with variance proxy $2\sigma_\varepsilon^2 = 2R_{\mathrm{obs}}$. Hence, by the sub-Gaussian tail bound and the noise gate~\eqref{eq:noise-gate},
\begin{align*}
  \mathbb{P}\bigl(\mathrm{Ared}_k < -\tfrac{\kappa_{\mathrm{tol}}}{2}\sqrt{2R_{\mathrm{obs}}} \,\big|\, \mathcal{F}_{k-1}\bigr)
  &\;\le\;
  \exp\Bigl(-\frac{\bigl(\tfrac{\kappa_{\mathrm{tol}}}{2}\sqrt{2R_{\mathrm{obs}}} + \tfrac{1}{4}\mathrm{Pred}_k\bigr)^2}{4R_{\mathrm{obs}}}\Bigr) \\
  &\;\le\;
  \exp\Bigl(-\frac{\tfrac{\kappa_{\mathrm{tol}}^2}{4}\,2R_{\mathrm{obs}}}{4R_{\mathrm{obs}}}\Bigr)
  \;=\; e^{-\kappa_{\mathrm{tol}}^2/8},
\end{align*}
where the last inequality uses $\mathrm{Pred}_k > 0$, and $p' \le \tfrac{1}{2}$ for $\kappa_{\mathrm{tol}} \ge \sqrt{8\ln 2}$ since $e^{-\ln 2} = \tfrac{1}{2}$.
\end{proof}

Lemma~\ref{lem:reject-prob} demonstrates that in the reliable regularization regime, trial step rejections are rare events with exponentially decaying probability $p' \le e^{-\kappa_{\mathrm{tol}}^2/8} \le \tfrac{1}{2}$. While this guarantees frequent successful steps in practice, establishing an almost-sure upper bound on $\sigma_k$ solely from this mechanism would require uniform geometric replenishment and lower bounds on $\mathbb{P}(r_k > \eta)$, which are difficult to verify a priori under arbitrary temporal drift. We therefore adopt regularization confinement as a standing assumption.

\begin{assumption}
\label{ass:sigma-confinement}
\textup{\textbf{(Regularization confinement.)}}
For the target accuracy $\varepsilon > 0$ with $\varepsilon \le 8L_g^2/\sigma_0$ and $\kappa_{\mathrm{tol}} \ge \sqrt{8\ln 2}$, the regularization parameter satisfies
\begin{equation}\label{eq:sigma-max}
  \sigma^*(\vct{g}_k) \;\le\; \sigma_k \;\le\; \sigma_{\max} \;:=\; \gamma_2\,\frac{8L_g^2}{\varepsilon},
\end{equation}
throughout all iterations $k$ prior to termination, i.e., while $\min_{j \le k} \|\nabla f(\vct{x}_j)\|_2 \ge \varepsilon$.
\end{assumption}

Complementing the probabilistic rejection bound of Lemma~\ref{lem:reject-prob}, the following proposition provides a deterministic upper bound on the number of unsuccessful iterations along any sample trajectory under regularization confinement.

\begin{proposition}
\label{prop:rejected-bound}
\textup{\textbf{(Bound on unsuccessful iterations.)}}
Let $\mathcal{S}_K := \{k \le K : \text{iteration } k \text{ is accepted}\}$ and $\mathcal{U}_K := \{k \le K : \text{iteration } k \text{ is rejected}\}$ denote the index sets of successful and unsuccessful iterations up to iteration $K$, respectively. Then:
\begin{equation}\label{eq:rejected-count}
  |\mathcal{U}_K| \;\le\; \frac{\log \gamma_1}{\log \gamma_2}\,|\mathcal{S}_K| \;+\; \frac{1}{\log \gamma_2}\log\frac{\sigma_{\max}}{\sigma_0}.
\end{equation}
\end{proposition}

\begin{proof}
From the update rule in Algorithm~\ref{alg:tobyqa}, the weight decreases by $\gamma_1$ only on accepted iterations with $r_k > \eta$, is held fixed on accepted iterations with $r_k \le \eta$, and increases by $\gamma_2$ on rejected iterations. Let $N_1 \le |\mathcal{S}_K|$ denote the number of accepted iterations on which $r_k > \eta$. Taking logarithms and telescoping from $k = 0$ to $K-1$ gives
\[
  \log \sigma_K \;=\; \log \sigma_0 \;-\; N_1\,\log \gamma_1 \;+\; |\mathcal{U}_K|\,\log \gamma_2.
\]
Since $N_1 \le |\mathcal{S}_K|$ and $\sigma_K \le \sigma_{\max}$ by Assumption~\ref{ass:sigma-confinement}, rearranging yields
\[
  |\mathcal{U}_K|\,\log \gamma_2
  \;=\; \log\frac{\sigma_K}{\sigma_0} + N_1 \log \gamma_1
  \;\le\; \log\frac{\sigma_{\max}}{\sigma_0} + |\mathcal{S}_K|\,\log \gamma_1,
\]
and dividing by $\log \gamma_2 > 0$ establishes~\eqref{eq:rejected-count}.
\end{proof}

\subsection{Iteration and Oracle Complexity}
\label{sec:expected-complexity}

We now state the iteration and oracle complexity bound for TOBYQA.

\begin{theorem}
\label{thm:main-complexity}
\textup{\textbf{(Oracle call complexity.)}}
Let Assumptions~\ref{ass:smoothness}--\ref{ass:sigma-confinement} hold. Let $\varepsilon > 0$ be any target gradient accuracy with $\varepsilon \le 8L_g^2/\sigma_0$ exceeding the statistical noise resolution floor
\begin{equation}\label{eq:noise-floor}
  \varepsilon \;\ge\; \varepsilon_{\mathrm{noise}} \;:=\; C_0\,\kappa_{eg}\,\max\Bigl\{\delta_{\mathrm{hi}},\; \frac{\sigma_\varepsilon}{\delta_{\mathrm{lo}}}\Bigr\},
\end{equation}
where $C_0 > 0$ is an absolute constant. Then the total number of iterations $K_\varepsilon$ (and consequently the number of oracle evaluations after initialization) required to achieve $\|\nabla f(\vct{x}_k)\|_2 \le \varepsilon$ satisfies
\begin{equation}\label{eq:complexity-rate}
  K_\varepsilon \;=\; O\Biggl(\frac{L_g\,\bigl(f(\vct{x}_0) - f_{\mathrm{low}}\bigr)}{\varepsilon^2}\Biggr).
\end{equation}
\end{theorem}

The restriction $\varepsilon\ge\varepsilon_{\mathrm{noise}}$ reflects the statistical resolution of single noisy queries: without replication, reductions below the observation-noise scale cannot be reliably distinguished from random fluctuations. The posterior-covariance stopping rule in Algorithm~\ref{alg:tobyqa} is calibrated to this same scale.

\begin{proof}
Let $k \in \mathcal{S}_K$ be an accepted iteration prior to termination, which implies that $\|\nabla f(\vct{x}_k)\|_2 \ge \varepsilon$. Under Assumption~\ref{ass:fully-linear} and condition~\eqref{eq:noise-floor}, the gradient estimation error satisfies
\[
  \|\vct{g}_k - \nabla f(\vct{x}_k)\|_2 \;\le\; \kappa_{eg}\,\zeta_k \;\le\; \frac{\varepsilon}{C_0} \;\le\; \frac{\|\nabla f(\vct{x}_k)\|_2}{C_0}.
\]
Choosing $C_0 \ge 8$ gives $\|\vct{g}_k - \nabla f(\vct{x}_k)\|_2 \le \tfrac{1}{8}\|\nabla f(\vct{x}_k)\|_2$, so by the reverse triangle inequality,
\[
  \|\vct{g}_k\|_2 \;\ge\; \|\nabla f(\vct{x}_k)\|_2 - \|\vct{g}_k - \nabla f(\vct{x}_k)\|_2 \;\ge\; \tfrac{7}{8}\|\nabla f(\vct{x}_k)\|_2 \;\ge\; \tfrac{1}{2}\,\varepsilon.
\]
Moreover, the model-error contribution satisfies
\begin{align*}
  \|\vct{g}_k - \nabla f(\vct{x}_k)\|_2\,\|\vct{s}_k\|_2
  &\;\le\; \frac{\varepsilon}{C_0}\,\|\vct{s}_k\|_2
  \;\le\; \frac{2\|\vct{g}_k\|_2}{C_0}\,\|\vct{s}_k\|_2 \\
  &\;=\; \frac{2}{C_0}\,\mathrm{Pred}_k
  \;\le\; \tfrac{1}{4}\,\mathrm{Pred}_k.
\end{align*}
Since $\sigma_k \ge \sigma^*(\vct{g}_k)$ by Assumption~\ref{ass:sigma-confinement}, Lemma~\ref{lem:reliable-sigma} guarantees that
\[
  f(\vct{x}_k) - f(\vct{x}_k + \vct{s}_k)
  \;\ge\; \tfrac{3}{4}\,\mathrm{Pred}_k - \|\vct{g}_k - \nabla f(\vct{x}_k)\|_2\,\|\vct{s}_k\|_2
  \;\ge\; \tfrac{1}{2}\,\mathrm{Pred}_k.
\]
On an accepted iteration $k \in \mathcal{S}_K$, the iterate updates to $\vct{x}_{k+1} = \vct{x}_k + \vct{s}_k$. Thus, using $\mathrm{Pred}_k = \|\vct{g}_k\|_2^{3/2} \sigma_k^{-1/2} \ge (\varepsilon/2)^{3/2} \sigma_{\max}^{-1/2}$ and the upper bound $\sigma_{\max} = \gamma_2 \frac{8L_g^2}{\varepsilon}$ from~\eqref{eq:sigma-max}, the objective reduction satisfies
\begin{equation}\label{eq:per-step-descent}
  f(\vct{x}_k) - f(\vct{x}_{k+1})
  \;\ge\;
  \tfrac{1}{4}\,\Bigl(\frac{\varepsilon}{2}\Bigr)^{3/2} \sqrt{\frac{\varepsilon}{8\gamma_2 L_g^2}}
  \;\ge\;
  c\,\frac{\varepsilon^2}{L_g},
\end{equation}
where $c = (32\sqrt{\gamma_2})^{-1} > 0$ is an absolute constant.

On rejected iterations $k \in \mathcal{U}_K$, the iterate remains unchanged ($\vct{x}_{k+1} = \vct{x}_k$), so $f(\vct{x}_k) - f(\vct{x}_{k+1}) = 0$. Summing the descent inequality~\eqref{eq:per-step-descent} over all accepted iterations up to iteration $K$ yields
\[
  f(\vct{x}_0) - f_{\mathrm{low}}
  \;\ge\; f(\vct{x}_0) - f(\vct{x}_K)
  \;=\; \sum_{k \in \mathcal{S}_K} \bigl(f(\vct{x}_k) - f(\vct{x}_{k+1})\bigr)
  \;\ge\; |\mathcal{S}_K|\,c\,\frac{\varepsilon^2}{L_g},
\]
whence the total number of accepted iterations satisfies
\begin{equation}\label{eq:accepted-bound}
  |\mathcal{S}_K| \;\le\; \frac{1}{c}\,\frac{L_g\,\bigl(f(\vct{x}_0) - f_{\mathrm{low}}\bigr)}{\varepsilon^2}
  \;=\; O\Biggl(\frac{L_g\,\bigl(f(\vct{x}_0) - f_{\mathrm{low}}\bigr)}{\varepsilon^2}\Biggr).
\end{equation}
By Proposition~\ref{prop:rejected-bound}, the number of unsuccessful iterations is bounded by
\[
  |\mathcal{U}_K| \;\le\; \frac{\log \gamma_1}{\log \gamma_2}\,|\mathcal{S}_K| \;+\; \frac{1}{\log \gamma_2}\log\frac{\sigma_{\max}}{\sigma_0}.
\]
Since $\sigma_{\max}/\sigma_0 = \gamma_2 (8L_g^2)/(\sigma_0\varepsilon) = O(\varepsilon^{-1})$, the additive logarithmic term satisfies
\[
  \frac{1}{\log \gamma_2}\log\frac{\sigma_{\max}}{\sigma_0} \;=\; O(\log(\varepsilon^{-1})),
\]
which is asymptotically dominated by $O(\varepsilon^{-2})$. Combining the bounds on successful and unsuccessful iterations gives
\[
  K_\varepsilon \;=\; |\mathcal{S}_K| + |\mathcal{U}_K| \;\le\; \Bigl(1 + \frac{\log \gamma_1}{\log \gamma_2}\Bigr)\,|\mathcal{S}_K| \;+\; O(\log(\varepsilon^{-1})) \;=\; O\Biggl(\frac{L_g\,\bigl(f(\vct{x}_0) - f_{\mathrm{low}}\bigr)}{\varepsilon^2}\Biggr).
\]
Since TOBYQA evaluates one new query per iteration after the initial $2n+1$ evaluations, the total number of oracle calls satisfies $N_{\mathrm{eval}} = 2n + 1 + K_\varepsilon = O(L_g(f(\vct{x}_0) - f_{\mathrm{low}})\varepsilon^{-2})$, completing the proof.
\end{proof}

\begin{remark}
\label{rem:optimality}
\textup{\textbf{(Comparison with first-order complexity.)}}
The rate $O(\varepsilon^{-2})$ in Theorem~\ref{thm:main-complexity} matches the information-theoretic lower bound for first-order methods on $C^{1,1}$ nonconvex objectives~\cite{CarmonDuchiHinderSidford2020}. This is consistent with the information available to TOBYQA at the budget $m=2n+1$: the active set supports a fully-linear model, but not an asymptotically accurate Hessian. The $O(\varepsilon^{-3/2})$ rate of derivative-based adaptive cubic regularization~\cite{NesterovPolyak2006,CartisGouldToint2011ARC,CartisGouldToint2011ARC2} requires $\|(H_k-\nabla^2f(\vct{x}_k))\vct{s}_k\|_2=O(\|\vct{s}_k\|_2^2)$~\cite[AM.4]{CartisGouldToint2011ARC}, whose recovery would generally require $O(n^2)$ samples. Under temporal drift, such a sample set also spans $O(n^2)$ time steps and incurs the nonlinear drift error quantified in Theorem~\ref{thm:nonlinear-drift}. TOBYQA instead uses an $O(n)$ active set, with a shorter temporal span and the first-order complexity rate.
\end{remark}


